% Options for packages loaded elsewhere
% Options for packages loaded elsewhere
\PassOptionsToPackage{unicode}{hyperref}
\PassOptionsToPackage{hyphens}{url}
%
\documentclass[
  english,
  russian,
  12pt,
  a4paper,
  DIV=11,
  numbers=noendperiod]{scrreprt}
\usepackage{xcolor}
\usepackage{amsmath,amssymb}
\setcounter{secnumdepth}{5}
\usepackage{iftex}
\ifPDFTeX
  \usepackage[T1]{fontenc}
  \usepackage[utf8]{inputenc}
  \usepackage{textcomp} % provide euro and other symbols
\else % if luatex or xetex
  \usepackage{unicode-math} % this also loads fontspec
  \defaultfontfeatures{Scale=MatchLowercase}
  \defaultfontfeatures[\rmfamily]{Ligatures=TeX,Scale=1}
\fi
\usepackage{lmodern}
\ifPDFTeX\else
  % xetex/luatex font selection
\fi
% Use upquote if available, for straight quotes in verbatim environments
\IfFileExists{upquote.sty}{\usepackage{upquote}}{}
\IfFileExists{microtype.sty}{% use microtype if available
  \usepackage[]{microtype}
  \UseMicrotypeSet[protrusion]{basicmath} % disable protrusion for tt fonts
}{}
\usepackage{setspace}
% Make \paragraph and \subparagraph free-standing
\makeatletter
\ifx\paragraph\undefined\else
  \let\oldparagraph\paragraph
  \renewcommand{\paragraph}{
    \@ifstar
      \xxxParagraphStar
      \xxxParagraphNoStar
  }
  \newcommand{\xxxParagraphStar}[1]{\oldparagraph*{#1}\mbox{}}
  \newcommand{\xxxParagraphNoStar}[1]{\oldparagraph{#1}\mbox{}}
\fi
\ifx\subparagraph\undefined\else
  \let\oldsubparagraph\subparagraph
  \renewcommand{\subparagraph}{
    \@ifstar
      \xxxSubParagraphStar
      \xxxSubParagraphNoStar
  }
  \newcommand{\xxxSubParagraphStar}[1]{\oldsubparagraph*{#1}\mbox{}}
  \newcommand{\xxxSubParagraphNoStar}[1]{\oldsubparagraph{#1}\mbox{}}
\fi
\makeatother


\usepackage{longtable,booktabs,array}
\newcounter{none} % for unnumbered tables
\usepackage{calc} % for calculating minipage widths
% Correct order of tables after \paragraph or \subparagraph
\usepackage{etoolbox}
\makeatletter
\patchcmd\longtable{\par}{\if@noskipsec\mbox{}\fi\par}{}{}
\makeatother
% Allow footnotes in longtable head/foot
\IfFileExists{footnotehyper.sty}{\usepackage{footnotehyper}}{\usepackage{footnote}}
\makesavenoteenv{longtable}
\usepackage{graphicx}
\makeatletter
\newsavebox\pandoc@box
\newcommand*\pandocbounded[1]{% scales image to fit in text height/width
  \sbox\pandoc@box{#1}%
  \Gscale@div\@tempa{\textheight}{\dimexpr\ht\pandoc@box+\dp\pandoc@box\relax}%
  \Gscale@div\@tempb{\linewidth}{\wd\pandoc@box}%
  \ifdim\@tempb\p@<\@tempa\p@\let\@tempa\@tempb\fi% select the smaller of both
  \ifdim\@tempa\p@<\p@\scalebox{\@tempa}{\usebox\pandoc@box}%
  \else\usebox{\pandoc@box}%
  \fi%
}
% Set default figure placement to htbp
\def\fps@figure{htbp}
\makeatother



\ifLuaTeX
\usepackage[bidi=basic,shorthands=off,provide=*]{babel}
\else
\usepackage[bidi=default,shorthands=off,provide=*]{babel}
\fi
\ifLuaTeX
  \usepackage{selnolig} % disable illegal ligatures
\fi


\setlength{\emergencystretch}{3em} % prevent overfull lines

\providecommand{\tightlist}{%
  \setlength{\itemsep}{0pt}\setlength{\parskip}{0pt}}



 
\usepackage[style=gost-numeric,backend=biber,langhook=extras,autolang=other*]{biblatex}
\addbibresource{bib/cite.bib}

\usepackage[]{csquotes}

\usepackage{indentfirst}
\usepackage{float}
\floatplacement{figure}{H}
\IfFileExists{plex-otf.sty}{
  %% Full TeXlive
  \usepackage[
    % math,
    RM={Scale=0.94},SS={Scale=0.94},SScon={Scale=0.94},TT={Scale=MatchLowercase,FakeStretch=0.9},
    DefaultFeatures={Ligatures=Common}
  ]{plex-otf}
  % \usepackage[Scale=MatchUppercase]{juliamono}
}{
  %% TinyTeX
  \usepackage{libertine}
}
\KOMAoption{captions}{tableheading}
\makeatletter
\@ifpackageloaded{caption}{}{\usepackage{caption}}
\AtBeginDocument{%
\ifdefined\contentsname
  \renewcommand*\contentsname{Содержание}
\else
  \newcommand\contentsname{Содержание}
\fi
\ifdefined\listfigurename
  \renewcommand*\listfigurename{Список иллюстраций}
\else
  \newcommand\listfigurename{Список иллюстраций}
\fi
\ifdefined\listtablename
  \renewcommand*\listtablename{Список таблиц}
\else
  \newcommand\listtablename{Список таблиц}
\fi
\ifdefined\figurename
  \renewcommand*\figurename{Рисунок}
\else
  \newcommand\figurename{Рисунок}
\fi
\ifdefined\tablename
  \renewcommand*\tablename{Таблица}
\else
  \newcommand\tablename{Таблица}
\fi
}
\@ifpackageloaded{float}{}{\usepackage{float}}
\floatstyle{ruled}
\@ifundefined{c@chapter}{\newfloat{codelisting}{h}{lop}}{\newfloat{codelisting}{h}{lop}[chapter]}
\floatname{codelisting}{Список}
\newcommand*\listoflistings{\listof{codelisting}{Листинги}}
\makeatother
\makeatletter
\makeatother
\makeatletter
\@ifpackageloaded{caption}{}{\usepackage{caption}}
\@ifpackageloaded{subcaption}{}{\usepackage{subcaption}}
\makeatother
\usepackage{bookmark}
\IfFileExists{xurl.sty}{\usepackage{xurl}}{} % add URL line breaks if available
\urlstyle{same}
% fallback for those not using the hyperref driver hyperxmp:
\makeatletter
\@ifundefined{xmpquote}{\newcommand{\xmpquote}[1]{#1}}{}
\makeatother
\hypersetup{
  pdftitle={Отчёт по лабораторной работе LAB04},
  pdfauthor={Mehmet Efe Kantoz},
  pdflang={ru-RU},
  hidelinks,
  pdfcreator={LaTeX via pandoc}}


\title{Отчёт по лабораторной работе LAB04}
\usepackage{etoolbox}
\makeatletter
\providecommand{\subtitle}[1]{% add subtitle to \maketitle
  \apptocmd{\@title}{\par {\large #1 \par}}{}{}
}
\makeatother
\subtitle{Модель гармонических колебаний}
\author{Mehmet Efe Kantoz}
\date{}
\begin{document}
\maketitle

\renewcommand*\contentsname{Содержание}
{
\setcounter{tocdepth}{1}
\tableofcontents
}
\listoffigures
\listoftables

\setstretch{1.5}
\section{Цель
работы}\label{ux446ux435ux43bux44c-ux440ux430ux431ux43eux442ux44b}

Изучить модель линейного гармонического осциллятора, освоить алгоритм
перехода от дифференциального уравнения второго порядка к системе двух
уравнений первого порядка, построить решения и фазовые портреты
гармонических колебаний для трех различных случаев (без затухания, с
затуханием, под действием внешней силы).

\section{Задание}\label{ux437ux430ux434ux430ux43dux438ux435}

\begin{enumerate}
\def\labelenumi{\arabic{enumi}.}
\tightlist
\item
  Построить решение уравнения гармонического осциллятора без затухания и
  его фазовый портрет.
\item
  Записать уравнение свободных колебаний гармонического осциллятора с
  затуханием, построить его решение и фазовый портрет.
\item
  Записать уравнение колебаний гармонического осциллятора при наличии
  внешней силы, построить его решение и фазовый портрет.
\item
  Ответить на контрольные вопросы к лабораторной работе.
\end{enumerate}

\chapter{Теоретическое
введение}\label{ux442ux435ux43eux440ux435ux442ux438ux447ux435ux441ux43aux43eux435-ux432ux432ux435ux434ux435ux43dux438ux435}

Движение грузика на пружине, математического маятника, колебания заряда
в электрическом контуре описываются одним и тем же дифференциальным
уравнением второго порядка:

\[
\ddot{x} + 2\gamma\dot{x} + \omega_0^2 x = f(t)
\]

где: - \(x\) --- переменная, описывающая состояние системы (смещение,
заряд); - \(\gamma\) --- параметр, характеризующий потери энергии
(трение, сопротивление); - \(\omega_0\) --- собственная частота
колебаний; - \(f(t)\) --- внешняя вынуждающая сила.

Для решения уравнения второго порядка выполняется переход к системе двух
дифференциальных уравнений первого порядка путем введения новой
переменной \(y = \dot{x}\):

\[
\begin{cases}
\dfrac{dx}{dt} = y \\
\dfrac{dy}{dt} = -2\gamma y - \omega_0^2 x + f(t)
\end{cases}
\]

Независимые переменные \(x\) и \(y\) образуют двумерное фазовое
пространство (фазовую плоскость). Состояние системы в каждый момент
времени задается точкой \((x(t), y(t))\), а эволюция системы
отображается в виде фазовой траектории.

\chapter{Выполнение лабораторной
работы}\label{ux432ux44bux43fux43eux43bux43dux435ux43dux438ux435-ux43bux430ux431ux43eux440ux430ux442ux43eux440ux43dux43eux439-ux440ux430ux431ux43eux442ux44b}

Зададим собственные параметры осциллятора и начальные условия: -
Собственная частота: \(\omega_0^2 = 3.5\) - Коэффициент затухания:
\(\gamma = 0.5\) - Начальное состояние: \(x(0) = 1.0\), \(y(0) = -0.5\)
- Внешняя сила: \(f(t) = 0.7 \sin(2t)\) - Интервал моделирования:
\(t \in [0, 50]\)

\section{Построение моделей в
Julia}\label{ux43fux43eux441ux442ux440ux43eux435ux43dux438ux435-ux43cux43eux434ux435ux43bux435ux439-ux432-julia}

Все три случая (консервативный осциллятор, свободные колебания с
затуханием, вынужденные колебания) программно реализованы в едином блоке
кода.

```julia using DifferentialEquations, Plots

\chapter{--- ОБЩИЕ ПАРАМЕТРЫ И НАЧАЛЬНЫЕ УСЛОВИЯ
---}\label{ux43eux431ux449ux438ux435-ux43fux430ux440ux430ux43cux435ux442ux440ux44b-ux438-ux43dux430ux447ux430ux43bux44cux43dux44bux435-ux443ux441ux43bux43eux432ux438ux44f}

w0\_sq = 3.5 \# w0\^{}2 (квадрат собственной частоты) gamma = 0.5 \# g
(параметр затухания) u0 = {[}1.0, -0.5{]} \# {[}x(0), x'(0){]} tspan =
(0.0, 50.0) \# Интервал времени t

\chapter{==============================================================================}\label{section}

\chapter{1. МОДЕЛЬ 1: Гармонический осциллятор без затухания (x'\,' +
w0\^{}2 * x =
0)}\label{ux43cux43eux434ux435ux43bux44c-1-ux433ux430ux440ux43cux43eux43dux438ux447ux435ux441ux43aux438ux439-ux43eux441ux446ux438ux43bux43bux44fux442ux43eux440-ux431ux435ux437-ux437ux430ux442ux443ux445ux430ux43dux438ux44f-x-w02-x-0}

\chapter{==============================================================================}\label{section-1}

function model1!(du, u, p, t) du{[}1{]} = u{[}2{]} du{[}2{]} = -w0\_sq *
u{[}1{]} end

prob1 = ODEProblem(model1!, u0, tspan) sol1 = solve(prob1, saveat=0.05)

\chapter{График зависимости
x(t)}\label{ux433ux440ux430ux444ux438ux43a-ux437ux430ux432ux438ux441ux438ux43cux43eux441ux442ux438-xt}

p1\_time = plot(sol1, vars=(0, 1), label=\enquote{x(t)},
xlabel=\enquote{Время t}, ylabel=\enquote{Координата x},
title=\enquote{1. Колебания без затухания}, color=:blue, lw=1.5)

\chapter{Фазовый портрет (x' от
x)}\label{ux444ux430ux437ux43eux432ux44bux439-ux43fux43eux440ux442ux440ux435ux442-x-ux43eux442-x}

p1\_phase = plot(sol1, vars=(1, 2), label=\enquote{Фазовая траектория},
xlabel=\enquote{Координата x}, ylabel=\enquote{Скорость x'},
title=\enquote{1. Фазовый портрет (без затухания)}, color=:darkblue,
lw=1.5)

\chapter{==============================================================================}\label{section-2}

\chapter{\texorpdfstring{2. МОДЕЛЬ 2: Свободные колебания с затуханием
(x'\,' + 2g\emph{x' + w0\^{}2 } x =
0)}{2. МОДЕЛЬ 2: Свободные колебания с затуханием (x'\,' + 2gx' + w0\^{}2  x = 0)}}\label{ux43cux43eux434ux435ux43bux44c-2-ux441ux432ux43eux431ux43eux434ux43dux44bux435-ux43aux43eux43bux435ux431ux430ux43dux438ux44f-ux441-ux437ux430ux442ux443ux445ux430ux43dux438ux435ux43c-x-2gx-w02-x-0}

\chapter{==============================================================================}\label{section-3}

function model2!(du, u, p, t) du{[}1{]} = u{[}2{]} du{[}2{]} = -2 *
gamma * u{[}2{]} - w0\_sq * u{[}1{]} end

prob2 = ODEProblem(model2!, u0, tspan) sol2 = solve(prob2, saveat=0.05)

\chapter{График зависимости
x(t)}\label{ux433ux440ux430ux444ux438ux43a-ux437ux430ux432ux438ux441ux438ux43cux43eux441ux442ux438-xt-1}

p2\_time = plot(sol2, vars=(0, 1), label=\enquote{x(t)},
xlabel=\enquote{Время t}, ylabel=\enquote{Координата x},
title=\enquote{2. Колебания с затуханием}, color=:red, lw=1.5)

\chapter{Фазовый
портрет}\label{ux444ux430ux437ux43eux432ux44bux439-ux43fux43eux440ux442ux440ux435ux442}

p2\_phase = plot(sol2, vars=(1, 2), label=\enquote{Фазовая траектория},
xlabel=\enquote{Координата x}, ylabel=\enquote{Скорость x'},
title=\enquote{2. Фазовый портрет (с затуханием)}, color=:maroon,
lw=1.5)

\chapter{==============================================================================}\label{section-4}

\chapter{\texorpdfstring{3. МОДЕЛЬ 3: Колебания под действием внешней
силы (x'\,' + 2g\emph{x' + w0\^{}2 } x =
f(t))}{3. МОДЕЛЬ 3: Колебания под действием внешней силы (x'\,' + 2gx' + w0\^{}2  x = f(t))}}\label{ux43cux43eux434ux435ux43bux44c-3-ux43aux43eux43bux435ux431ux430ux43dux438ux44f-ux43fux43eux434-ux434ux435ux439ux441ux442ux432ux438ux435ux43c-ux432ux43dux435ux448ux43dux435ux439-ux441ux438ux43bux44b-x-2gx-w02-x-ft}

\chapter{==============================================================================}\label{section-5}

f(t) = 0.7 * sin(2.0 * t)

function model3!(du, u, p, t) du{[}1{]} = u{[}2{]} du{[}2{]} = -2 *
gamma * u{[}2{]} - w0\_sq * u{[}1{]} + f(t) end

prob3 = ODEProblem(model3!, u0, tspan) sol3 = solve(prob3, saveat=0.05)

\chapter{График зависимости
x(t)}\label{ux433ux440ux430ux444ux438ux43a-ux437ux430ux432ux438ux441ux438ux43cux43eux441ux442ux438-xt-2}

p3\_time = plot(sol3, vars=(0, 1), label=\enquote{x(t)},
xlabel=\enquote{Время t}, ylabel=\enquote{Координата x},
title=\enquote{3. Колебания с внешней силой}, color=:green, lw=1.5)

\chapter{Фазовый
портрет}\label{ux444ux430ux437ux43eux432ux44bux439-ux43fux43eux440ux442ux440ux435ux442-1}

p3\_phase = plot(sol3, vars=(1, 2), label=\enquote{Фазовая траектория},
xlabel=\enquote{Координата x}, ylabel=\enquote{Скорость x'},
title=\enquote{3. Фазовый портрет (с внешней силой)}, color=:darkgreen,
lw=1.5)

\chapter{--- СОХРАНЕНИЕ ГРАФИКОВДЛЯ ОТЧЕТА
---}\label{ux441ux43eux445ux440ux430ux43dux435ux43dux438ux435-ux433ux440ux430ux444ux438ux43aux43eux432ux434ux43bux44f-ux43eux442ux447ux435ux442ux430}

grafik 1(рис.~\ref{fig-001})

\begin{figure}

\centering{

\includegraphics[width=0.7\linewidth,height=\textheight,keepaspectratio]{image/1.png}

}

\caption{\label{fig-001}grafik1}

\end{figure}%

grafik 2(рис.~\ref{fig-002})

\begin{figure}

\centering{

\includegraphics[width=0.7\linewidth,height=\textheight,keepaspectratio]{image/2.png}

}

\caption{\label{fig-002}grafik2}

\end{figure}%

grafik 3(рис.~\ref{fig-003})

\begin{figure}

\centering{

\includegraphics[width=0.7\linewidth,height=\textheight,keepaspectratio]{image/3.png}

}

\caption{\label{fig-003}grafik3}

\end{figure}%

grafik 4(рис.~\ref{fig-004})

\begin{figure}

\centering{

\includegraphics[width=0.7\linewidth,height=\textheight,keepaspectratio]{image/4.png}

}

\caption{\label{fig-004}grafik4}

\end{figure}%

grafik 5(\textbf{?@fig-005})

\begin{figure}

\centering{

\includegraphics[width=0.7\linewidth,height=\textheight,keepaspectratio]{image/5.png}

}

\caption{\label{fig-006}grafik5}

\end{figure}%


\printbibliography



\end{document}
